\uuid{AT3h}
\exo7id{4296}
\titre{exo7 4296}
\auteur{quercia}
\organisation{exo7}
\datecreate{2010-03-12}
\isIndication{false}
\isCorrection{true}
\chapitre{Calcul d'intégrales}
\sousChapitre{Intégrale impropre}
\cpge{mpsi-mp2i.integration.c, pcsi.integration.b, ptsi.integration.b, mpsi-mp2i.limites-continuite.a, pcsi.limites-continuite.a, ptsi.limites-continuite.a, bcpst1.limites-continuite.a, bcpst1.integration-segment.a, mp-mpi.integration-intervalle.b, bcpst2.integrales-generalisees}
\module{Analyse}
\niveau{L1}
\difficulte{}

\contenu{
\begin{enumerate}
    \item \question{Soit $f : {\R^+} \to {\R^+}$ continue ayant une limite finie en~$+\infty$.
    Montrer que $ \int_{t=0}^{+\infty} (f(t+a) - f(t))\,d t$ converge.}
    \item \question{Calculer $ \int_{t=0}^{+\infty} (\arctan(t+1) - \arctan(t))\,d t$.}
\reponse{
$\frac\pi4+\frac{\ln 2}2$.
}
\end{enumerate}
}
